%%%%%%%%%%%%%%%%%%%%%%%%%%%%%%%%%%%%%%%%%%%%%%%%%%%%%%%%%%%%%%%%%%%%%%%%%%%%%%%
%                         CAPITULO 4
%           Introduccion a las Redes Industriales
%%%%%%%%%%%%%%%%%%%%%%%%%%%%%%%%%%%%%%%%%%%%%%%%%%%%%%%%%%%%%%%%%%%%%%%%%%%%%%%

\chapter{Introduccion a las Redes Industriales}

Hasta ahora hemos estudiado como comunican dispositivos mediante cables, ondas
de radio, y como se organizan las redes de computadores mediante direcciones IP,
switches y routers. Sin embargo, en un entorno de fabricacion o produccion, las
redes no solo conectan ordenadores: conectan sensores, actuadores, robots, PLCs
y sistemas de supervision, todo trabajando en tiempo real para controlar procesos
industriales.

En este capitulo introducimos el concepto de \textbf{redes industriales}, tambien
conocidas como redes de comunicacion para la automatizacion y control de
procesos. Veremos que son, como funcionan, que tipos existen y por que son
fundamentales en la industria moderna.


\section{Sistemas industriales distribuidos}

\subsection{Definicion}

Un \textbf{sistema industrial distribuido} es una arquitectura de control y
automatizacion donde multiples dispositivos, sensores, actuadores y sistemas
informaticos estan interconectados para operar de manera descentralizada.

\begin{definicion}{Sistema industrial distribuido}
Un \textbf{sistema industrial distribuido} es una arquitectura de control donde
multiples dispositivos (PLCs, sensores, actuadores, robots) comparten la
responsabilidad de la operacion, interconectados mediante una red de
comunicacion industrial. Ofrece mayor flexibilidad, escalabilidad y eficiencia
que los sistemas centralizados tradicionales.
\end{definicion}

En un sistema centralizado clasico, un unico controlador (PLC o computador)
toma todas las decisiones y envia ordenes a los dispositivos de campo. En un
sistema distribuido, la inteligencia se reparte: cada dispositivo puede tomar
decisiones locales y compartir informacion con el resto.

\subsection{Caracteristicas principales}

Los sistemas industriales distribuidos se caracterizan por:

\begin{itemize}
    \item \textbf{Descentralizacion}: multiples dispositivos comparten la
    responsabilidad de la operacion. No hay un unico punto de control.
    \item \textbf{Interconectividad}: se basan en redes de comunicacion industrial
    como Ethernet/IP, Modbus, Profibus o IoT industrial.
    \item \textbf{Escalabilidad}: se pueden agregar nuevos dispositivos sin afectar
    la operacion global.
    \item \textbf{Tolerancia a fallos}: al no depender de un unico nodo, el sistema
    puede seguir funcionando ante fallos parciales.
    \item \textbf{Optimizacion de recursos}: permite un uso eficiente de energia,
    maquinaria y mano de obra.
    \item \textbf{Automatizacion avanzada}: utiliza PLCs, SCADA, IIoT e inteligencia
    artificial.
\end{itemize}

\subsection{Aplicaciones y ventajas}

Las aplicaciones mas comunes incluyen:
\begin{itemize}
    \item \textbf{Automatizacion de fabricas}: produccion y ensamblaje de productos.
    \item \textbf{Redes de energia inteligente (Smart Grid)}: gestion eficiente de
    energia electrica.
    \item \textbf{Transporte y logistica}: control distribuido en almacenes y vehiculos
    autonomos.
    \item \textbf{Industria 4.0}: implementacion de IoT y Big Data para optimizacion.
    \item \textbf{Robotica industrial}: control distribuido en lineas de produccion.
\end{itemize}

Las ventajas principales son: mayor eficiencia operativa, reduccion de costes de
mantenimiento, mayor adaptabilidad a cambios en la produccion, y mejora en la
recopilacion y analisis de datos.


\section{Redes industriales: conceptos fundamentales}

\subsection{Definicion}

\begin{definicion}{Red industrial}
Una \textbf{red industrial} es una estructura de comunicacion automatizada que
gestiona procesos industriales, involucrando actuadores, computadoras, maquinas,
sensores, interfaces, medios de comunicacion (fibra optica, cables, redes
inalambricas) y robots. Su objetivo es transmitir y compartir datos para garantizar
un funcionamiento eficiente de los procesos productivos.
\end{definicion}

\subsection{Como funcionan}

Las redes industriales reemplazan el modelo tradicional de control cableado a
mano (donde cada dispositivo necesitaba su propio cable de control) por un modelo
de comunicacion digital donde multiples dispositivos comparten un unico medio de
transmision.

Este modelo reduce drasticamente la cantidad de cableado necesario, simplifica
la organizacion del sistema y permite una flexibilidad mucho mayor para
reconfigurar la produccion.

\subsection{Convergencia de TI y TO (OT)}

En la industria moderna se distinguen dos tecnologias:

\begin{itemize}
    \item \textbf{TI (Tecnologia de la Informacion)}: se ocupa de la entrega de
    informacion cuando un profesional la solicita. Sistemas de datos, ERP, CRM.
    \item \textbf{TO (Tecnologia Operacional) o OT}: trabaja con la entrega de
    informacion en un momento especifico bajo una condicion determinada.
    Controla equipos industriales en tiempo real.
\end{itemize}

La convergencia de TI y TO esta relacionada con el auge del \textbf{edge computing}
(computacion en el borde): trasladar los recursos informaticos a la ubicacion
fisica de la fuente de datos (por ejemplo, en la propia fabrica), permitiendo unificar
sistemas que tradicionalmente estaban aislados.

\subsection{Desafios}

Las redes industriales enfrentan desafios especificos:
\begin{itemize}
    \item \textbf{Complejidad en la integracion}: sistemas heterogeneos de diferentes
    fabricantes deben comunicarse.
    \item \textbf{Ciberseguridad}: las redes industriales son objetivos criticos; un ataque
    puede parar la produccion.
    \item \textbf{Capacitacion especializada}: requiere personal con conocimientos en
    automatizacion y redes.
\end{itemize}


\section{Clasificacion de las redes industriales}

Las redes industriales se clasifican segun el nivel de la piramide de automatizacion:

\begin{table}[htbp]
\centering
\small
\caption{Niveles de las redes industriales}
\begin{tabularx}{\textwidth}{@{}l>{\raggedright\arraybackslash}X>{\raggedright\arraybackslash}X@{}}
\toprule
\textbf{Nivel} & \textbf{Funcion} & \textbf{Ejemplos} \\
\midrule
Nivel de informacion (gestion) & Planificacion, ERP, MES & Redes Ethernet, TCP/IP, bases de datos \\
Nivel de control & Controladores, PLCs, SCADA & ControlNet, Ethernet/IP \\
Nivel de campo y proceso & Sensores, actuadores, transmisores & Profibus, Modbus, CAN \\
Nivel de entrada/salida & Dispositivos simples, fotocelulas & AS-i, Sensor Bus \\
\bottomrule
\end{tabularx}
\end{table}

\subsection{Clasificacion por funcionalidad}

Segun sus capacidades y el tipo de datos que transportan:

\begin{enumerate}
    \item \textbf{Sensor Bus (buses de sensores)}: conectan sensores digitales y
    actuadores simples. Transmite pocos datos, distancias cortas, bajo costo.
    Ejemplos: AS-i, CAN, Interbus.
    \item \textbf{Device Bus (buses de dispositivos)}: intermedian entre sensores y
    PLCs. Manejan variables de proceso (presion, temperatura, flujo).
    Alcance hasta 500 metros. Ejemplos: Modbus, Profibus DP, DeviceNet.
    \item \textbf{Field Bus (buses de campo)}: permiten que varios instrumentos se
    comuniquen en red para control y monitoreo. Conectan mas equipos a mayor
    distancia. Ejemplos: Foundation Fieldbus, Profibus PA.
    \item \textbf{Control Bus / Ethernet Industrial}: redes de alta velocidad para
    control en tiempo real. Ejemplos: Ethernet/IP, PROFINET, EtherCAT, ControlNet.
\end{enumerate}

\subsection{Tabla comparativa de tipos de redes industriales}

\begin{table}[htbp]
\centering
\small
\caption{Comparacion de tipos de redes industriales}
\begin{tabularx}{\textwidth}{@{}l>{\raggedright\arraybackslash}X>{\raggedright\arraybackslash}X>{\raggedright\arraybackslash}X@{}}
\toprule
\textbf{Tipo} & \textbf{Velocidad} & \textbf{Funcionalidad} & \textbf{Uso tipico} \\
\midrule
Sensor Bus & Baja & Simple & Sensores, actuadores binarios \\
Device Bus & Media & Media & Transmisores, servomotores \\
Field Bus & Media & Alta & Control de procesos continuos \\
Ethernet Industrial & Alta & Muy alta & Control en tiempo real, integracion TI/TO \\
\bottomrule
\end{tabularx}
\end{table}


\section{Principales protocolos de redes industriales}

\subsection{PROFINET}

PROFINET es una solucion de Ethernet Industrial abierta basada en estandares
internacionales. Es el protocolo de comunicacion desarrollado para intercambiar
datos entre controladores y dispositivos en un sector automatizado.

Utiliza Ethernet estandar como medio de comunicacion, lo que permite que otros
protocolos Ethernet (SNMP, MQTT, HTTP) coexistan en la misma infraestructura.
Ademas, define la comunicacion ciclica y aciclica entre componentes, incluyendo
diagnosticos, seguridad funcional, alarmas e informacion adicional.

\subsection{PROFIBUS}

PROFIBUS es una red digital que proporciona comunicacion entre sensores de
campo y controladores. Surgio inicialmente como PROFIBUS FMS (Fieldbus Message
Specification), pero hoy en dia los mas utilizados son:

\begin{itemize}
    \item \textbf{PROFIBUS DP} (Decentralised Peripherals): para automatizacion de
    fabricas, comunicacion rapida entre PLCs y dispositivos de campo.
    \item \textbf{PROFIBUS PA} (Process Automation): para procesos continuos como
    industrias petroquimica y minera.
\end{itemize}

\subsection{EtherCAT}

EtherCAT (Ethernet for Control Automation Technology) es un protocolo que lleva
la potencia y flexibilidad de Ethernet a la automatiz industrial, al control de
movimientos y a los sistemas de control en tiempo real.

Fue desarrollado por Beckhoff Automation y estandarizado bajo IEC 61158. Es
mantenido por el EtherCAT Technology Group (ETG). Su principal ventaja es la
extremada velocidad de transmision y la sincronizacion precisa de los dispositivos.

\subsection{ControlNet}

ControlNet utiliza el Protocolo Industrial Comun (CIP) para las capas superiores
del modelo OSI. Se desarrollo para ofrecer control confiable y de alta velocidad,
con transferencia de datos I/O programada en momentos especificos.

Garantiza que los mensajes criticos que no dependen del tiempo se ejecuten sin
interferir en la transferencia de datos de control. Se utiliza normalmente para
aplicaciones redundantes.

\subsection{Ethernet en la industria: importancia en la Industria 4.0}

Ethernet se ha convertido en el estandar mundial para la intercomunicacion de
datos en red porque:
\begin{itemize}
    \item Tiene una implementacion simple y costos reducidos.
    \item Permite el uso de varios protocolos en un entorno industrial.
    \item Es interoperable y escalable para las demandas.
    \item Sigue en constante evolucion y mejora.
\end{itemize}

En la \textbf{Industria 4.0}, la informacion fluye horizontal y verticalmente a traves
de toda la organizacion. Esto requiere una infraestructura de comunicacion que
abarque todas las maquinas, computadoras, robots y equipos, lo cual solo es
posible con Ethernet e Internet unidos al IIoT (Internet Industrial de las Cosas).


\section{Resumen de la Unidad 5}

\begin{itemize}
    \item Los \textbf{sistemas industriales distribuidos} descentralizan el control,
    permitiendo que multiples dispositivos compartan la responsabilidad de la
    operacion y mejoren la tolerancia a fallos.
    \item Una \textbf{red industrial} es una estructura de comunicacion que gestiona
    procesos industriales, conectando sensores, actuadores, PLCs, robots y
    sistemas de supervision.
    \item La \textbf{convergencia de TI y TO} (OT) unifica los sistemas de informacion
    empresarial con los sistemas de control en tiempo real, facilitando el
    \textbf{edge computing} y la Industria 4.0.
    \item Las redes industriales se clasifican en niveles: informacion, control,
    campo y E/S. Por funcionalidad: Sensor Bus, Device Bus, Field Bus y
    Ethernet Industrial.
    \item Los principales protocolos son: \textbf{PROFINET}, \textbf{PROFIBUS},
    \textbf{EtherCAT}, \textbf{ControlNet}, \textbf{Modbus}, \textbf{Ethernet/IP}.
    \item \textbf{Ethernet} es el estandar dominante en la industria moderna por su
    simplicidad, bajo costo, interoperabilidad y capacidad para soportar
    multiples protocolos simultaneamente.
\end{itemize}
